Das vorangegangene Kapitel hat die fachlichen und technischen Anforderungen an den CrawlDataTransformer hergeleitet und in einen Architekturentwurf überführt. Aufbauend darauf 
beschreibt dieses Kapitel die prototypische Umsetzung der zentralen Entwurfsentscheidungen. Im Mittelpunkt stehen dabei die konkreten Mechanismen, durch die unstrukturierte 
Produktdaten vorverarbeitet, extrahiert, geprüft, normalisiert und für die weitere Verarbeitung bereitgestellt werden.

Eine vollständige Beschreibung sämtlicher Klassen, Adapter und Persistenzkomponenten ist für die Zielsetzung der Arbeit nicht erforderlich. Die Darstellung konzentriert sich daher 
auf diejenigen Implementierungsaspekte, die für die Untersuchung der Extraktionsansätze, ihre Vergleichbarkeit und die anschließende Evaluation maßgeblich sind. Hierzu zählen 
insbesondere die gemeinsame HTML-Vorverarbeitung, die mehrphasige Transformationsverarbeitung, die LLM-basierte Extraktion einschließlich der Promptstrategie, der regelbasierte 
Referenzansatz für das Datenfeld \textit{ProductSizing} sowie die Schema-Validierung, Normalisierung und Erfassung der Evaluationsdaten.

Abschnitt 5.1 beschreibt zunächst den technologischen Implementierungsrahmen. Abschnitt 5.2 erläutert die technische Umsetzung der Transformationsverarbeitung und der gemeinsamen 
Vorverarbeitung. Die LLM-basierte Extraktion und ihre Promptstrategie werden in Abschnitt 5.3 dargestellt. Abschnitt 5.4 behandelt die Umsetzung des kaskadierten regelbasierten 
Referenzansatzes. Abschließend fasst Abschnitt 5.5 die nachgelagerte Validierung, die technische Unterstützung der Evaluation und die Anbindung an das Zielsystem zusammen.

\subsection{Technologischer Implementierungsrahmen}

Der technologische Implementierungsrahmen ergibt sich aus den in Kapitel 4 beschriebenen Anforderungen und den bestehenden Randbedingungen des Projektkontexts. Python, die Gemini API 
und MongoDB waren als grundlegende Technologien vorgegeben. Die ergänzend eingesetzten Frameworks und Bibliotheken wurden so gewählt, dass sie die typisierte Verarbeitung 
verschachtelter Produktdaten, die zeitlich getrennte Kommunikation mit externen Diensten, die schrittweise Persistierung von Zwischenständen sowie unterschiedliche Eingabe- und 
Ausgabewege unterstützen. Der Prototyp verwendet Python 3.12 als einheitliche Laufzeitbasis.

Die zentralen Technologien und ihre jeweilige Funktion innerhalb des CrawlDataTransformer sind in Tabelle~\ref{tab:technologien-implementierung} zusammengefasst.

\begin{table}[H]
    \centering
    \begingroup
    \small
    \setlength{\tabcolsep}{5pt}
    \renewcommand{\arraystretch}{1.18}
    \begin{tabularx}{\textwidth}{
        @{}
        >{\raggedright\arraybackslash}p{0.18\textwidth}
        >{\raggedright\arraybackslash}p{0.25\textwidth}
        >{\raggedright\arraybackslash}X
        @{}
    }
        \toprule
        \textbf{Bereich}
        & \textbf{Technologien}
        & \textbf{Funktion im Prototyp} \\
        \midrule

        Laufzeit und Modelle
        & Python 3.12, Pydantic
        & Implementierung des Anwendungskerns sowie typisierte Abbildung und Prüfung von Domänen-, Ergebnis- und Konfigurationsmodellen. \\
        \addlinespace[2pt]

        Nebenläufigkeit
        & \texttt{asyncio}
        & Asynchrone Bearbeitung von Ein- und Ausgaben, Begrenzung gleichzeitig verarbeiteter Anfragen und Koordination der beiden Extraktionszweige. \\
        \addlinespace[2pt]

        Persistenz
        & MongoDB, Beanie, Motor
        & Dokumentorientierte Speicherung verschachtelter und während der Verarbeitung schrittweise ergänzter Arbeits-, Zustands- und Evaluationsdaten. \\
        \addlinespace[2pt]

        Ein- und Ausgabe
        & FastAPI, Uvicorn, Click, HTTPX
        & Bereitstellung der REST- und Kommandozeilenschnittstelle sowie HTTP-basierte Kommunikation mit der Ziel-API. \\
        \addlinespace[2pt]

        HTML und LLM
        & Beautiful Soup, Google Gen AI SDK
        & Strukturorientierte Verarbeitung der HTML-Dokumente und technische Anbindung der Gemini API. \\
        \addlinespace[2pt]

        Ausführung
        & Docker Compose, PDM
        & Bereitstellung einer einheitlichen lokalen Laufzeitumgebung sowie Verwaltung und Festschreibung der Python-Abhängigkeiten. \\

        \bottomrule
    \end{tabularx}
    \caption{Zentrale Technologien des CrawlDataTransformer}
    \label{tab:technologien-implementierung}
    \endgroup
\end{table}

Python bildet die gemeinsame Grundlage für den Anwendungskern, die fachlichen Filter und die technischen Adapter. Die internen Datenstrukturen werden überwiegend als Pydantic-Modelle 
umgesetzt. Dadurch können Eingaben, Batchkonfigurationen, Verarbeitungsergebnisse und verschachtelte Produktobjekte typisiert abgebildet und an den Grenzen einzelner 
Verarbeitungsschritte geprüft werden. Die Modelle dienen zugleich als gemeinsame Repräsentation zwischen dem Anwendungskern und den technischen Adaptern. Datenbankspezifische 
Dokumentmodelle bleiben davon getrennt und werden über Mapper in die Domänenmodelle beziehungsweise aus diesen überführt.

Die Verarbeitung enthält mehrere I/O-intensive Operationen, insbesondere Zugriffe auf MongoDB sowie die Kommunikation mit der Gemini API und der Ziel-API. Hierfür verwendet der 
Prototyp die asynchrone Laufzeitunterstützung von \texttt{asyncio}. Diese ermöglicht die nebenläufige Bearbeitung mehrerer Transformationsanfragen und die Koordination des 
LLM-basierten und des regelbasierten Extraktionszweigs, ohne die Verarbeitung als parallele Berechnung über mehrere Prozessorkerne auszulegen. Die Anzahl gleichzeitig 
bearbeiteter Anfragen wird über die Batchkonfiguration begrenzt.

MongoDB wird als lokale Arbeits-, Transformations- und Evaluationsdatenbank eingesetzt. Das dokumentorientierte Datenmodell unterstützt die Speicherung der verschachtelten 
Verarbeitungskontexte, die im Verlauf der Pipeline um Vorverarbeitungs-, Extraktions-, Validierungs-, Normalisierungs- und Bewertungsergebnisse ergänzt werden. Beanie übernimmt 
die Abbildung der Dokumentmodelle, während Motor den asynchronen Datenbankzugriff bereitstellt. Die lokale Persistenz bleibt von der produktiven Datenhaltung getrennt. Transformierte 
Produktdaten werden ausschließlich über die Ziel-API übergeben.

Für die Eingabe stehen zwei technische Adapter zur Verfügung. FastAPI und Uvicorn realisieren die REST-Schnittstelle, während Click die kommandozeilenbasierte Ausführung mit 
dateibasierten Eingaben ermöglicht. Beide Adapter überführen ihre Eingaben in dieselben internen Modelle und greifen über den gemeinsamen Eingangsport auf dieselben Anwendungsfälle zu. 
Beautiful Soup wird für das Parsen und die strukturorientierte Bereinigung der HTML-Dokumente eingesetzt. Die Kommunikation mit Gemini erfolgt über das Google Gen AI SDK, während 
HTTPX die asynchrone Anbindung der Ziel-API übernimmt.

Die Quellcodestruktur bildet die in Kapitel 4 beschriebene Trennung der Verantwortlichkeiten ab. Der Bereich \texttt{core} enthält die Domänenmodelle, Filter, Pipelines, Prompts 
und anwendungsinternen Services. Die unter \texttt{ports} definierten Schnittstellen beschreiben die vom Anwendungskern benötigten Ein- und Ausgangsfunktionen. Ihre konkreten 
technischen Umsetzungen befinden sich unter \texttt{adapters}. Umgebungsabhängige Einstellungen werden im Konfigurationsbereich gekapselt. Die Zusammensetzung der Repositories, 
Adapter, Filter, Pipelines und Services erfolgt zentral im Dependency Wiring, sodass die fachlichen Komponenten nicht selbst für die Erzeugung ihrer technischen Abhängigkeiten 
verantwortlich sind.

\subsection{Umsetzung der Transformationsverarbeitung}

Die in Abschnitt 4.2.4 dargestellte zeitlich getrennte Verarbeitung wird durch mehrere request- und batchbezogene Teilpipelines umgesetzt. Der \texttt{TransformationService} 
koordiniert deren Ausführung, verwaltet die persistenten Arbeitsobjekte und führt die Ergebnisse der getrennten Extraktionszweige zusammen. Die konkrete Verarbeitung gliedert 
sich in die Vorbereitung einzelner Transformationsanfragen, die batchweite Einreichung der LLM-Anfragen, die regelbasierte Extraktion sowie die nachgelagerte Verarbeitung der 
abgerufenen LLM-Ergebnisse. Eine einzelne durchgängige und synchrone Filterkette wird damit nicht gebildet.

\subsubsection{Mehrphasige Pipelineausführung}

Zu Beginn eines Transformationslaufs erzeugt der \texttt{TransformationService} zunächst einen \texttt{BatchRequest} im Zustand \texttt{CREATED} und speichert diesen in MongoDB. 
Anschließend wird für jedes Eingabeobjekt ein eigener \texttt{TransformationRequest} angelegt und über seine Batchkennung dem übergeordneten Verarbeitungslauf zugeordnet. Erst 
nachdem die internen Kennungen und Beziehungen persistiert wurden, wechselt der Batch in den Zustand \texttt{RUNNING} und die fachliche Verarbeitung beginnt. Dadurch sind die 
ursprünglichen Eingaben bereits vor dem ersten Verarbeitungsschritt eindeutig adressierbar.

Die requestbezogene Verarbeitung wird durch den \texttt{TransformationPipelineExecutor} ausgeführt. Dieser überführt die zu bearbeitenden Anfragen in eine \texttt{asyncio.Queue} und 
erzeugt eine konfigurierbare Anzahl asynchroner Worker. Die Anzahl der Worker wird durch \texttt{max\_workers} begrenzt und zusätzlich an die Anzahl der tatsächlich vorhandenen 
Transformationsanfragen angepasst. Innerhalb einer einzelnen Anfrage werden die Filter sequenziell ausgeführt, sodass das Ergebnis eines Filters als Eingabe des folgenden Schritts 
dient. In den Vorbereitungs- und Nachverarbeitungspipelines wird der aktualisierte \texttt{TransformationRequest} nach jedem abgeschlossenen Filter persistiert. Damit bleiben bereits 
erzeugte Zwischenstände auch dann erhalten, wenn ein späterer Verarbeitungsschritt fehlschlägt.

Die initiale requestbezogene Pipeline enthält das HTML-Cleanup. Nach dessen Abschluss werden die LLM-Einreichung und die regelbasierte Extraktion als getrennte \texttt{asyncio}-Tasks 
gestartet. Hierfür erhalten beide Zweige tiefe Kopien derselben vorverarbeiteten Transformationsanfragen. Die LLM-Einreichung wird als batchbezogene Pipeline ausgeführt, da die 
geeigneten Anfragen gemeinsam an den externen Dienst übergeben werden. Der während dieser Ausführung erzeugte Batchkontext stellt unter anderem die Batchkennung und einen Snapshot der 
verwendeten Konfiguration bereit. Dazu gehören das Modell, die Promptversion, die Temperatur und das Thinking-Budget.

Die regelbasierte Extraktion arbeitet dagegen requestbezogen. Da diese Teilpipeline nur einen Extraktionsfilter enthält, wird sie für die einzelnen Anfragen direkt asynchron ausgeführt. 
Nach Abschluss beider Zweige werden das regelbasierte Ergebnis und der zugehörige Pipelinekontext anhand der internen Request-ID in die vom LLM-Zweig zurückgegebenen Anfragen 
übernommen. Anschließend werden die zusammengeführten Transformationsanfragen und die externe LLM-Batchreferenz persistiert. Die Ausführung ist damit fachlich parallel und asynchron 
koordiniert, stellt jedoch keine parallele Berechnung über mehrere Prozessorkerne dar.

Die LLM-Einreichung erzeugt zu diesem Zeitpunkt noch kein fachliches Produktergebnis. Die weitere Verarbeitung wird durch einen separaten Aufruf für den bereits persistierten 
Batch gestartet. Dabei prüft der \texttt{LlmBatchStatusRefreshService} zunächst den Zustand des externen LLM-Batches. Liegen die Antworten vor, werden sie anhand der übermittelten 
Request-IDs den ursprünglichen Transformationsanfragen zugeordnet und als Extraktionsergebnisse gespeichert. Nur Anfragen im Zustand \texttt{EXTRACTED} werden anschließend durch 
die Nachverarbeitungspipeline geführt. Diese umfasst Schema-Validierung, Normalisierung, Produkt-Matching, Evaluation und Export. Die Umsetzung bildet somit die zeitliche 
Trennung zwischen Einreichung und Ergebnisverarbeitung unmittelbar im Ablauf und in den persistenten Zuständen ab.

\subsubsection{Gemeinsame HTML-Vorverarbeitung}

Das HTML-Cleanup erzeugt die gemeinsame Eingabebasis für den LLM-basierten und den regelbasierten Extraktionszweig. Der zugehörige Filter verfolgt eine konservative 
Bereinigungsstrategie: Technische und mit hinreichender Sicherheit irrelevante Inhalte werden entfernt, während potenziell produktbezogene Informationen und strukturelle 
Hinweise erhalten bleiben. Eine fachliche Extraktion einzelner Produktattribute findet in diesem Schritt noch nicht statt.

Zu Beginn werden seitenweite Metadaten ausgelesen. Hierzu gehören der Seitentitel, die kanonische URL, der Open-Graph-Titel, der Open-Graph-Seitenname und die Sprache des 
HTML-Dokuments. Anschließend entfernt der Filter Kommentare sowie technische Elemente wie Skripte, Styles und eingebettete Medien. Weitere Bereiche werden anhand ihrer 
HTML-Tags, IDs und Klassen nur dann verworfen, wenn sie eindeutig einer globalen Navigation, dem Footer, Cookie- und Consent-Dialogen, Such- und Warenkorbfunktionen, 
Trackingelementen oder nicht produktbezogenen Empfehlungsbereichen zugeordnet werden können. Produktnahe oder uneindeutige Inhalte bleiben dagegen erhalten.

Bild-Tags werden aus dem HTML entfernt, da die Bilddaten bereits unabhängig davon Bestandteil der Eingabe sein können. Vorhandene Alt-Texte werden jedoch als Textknoten 
erhalten, da sie fachlich relevante Bezeichnungen enthalten können. Danach reduziert der Filter die HTML-Attribute auf strukturell oder semantisch nutzbare Informationen. 
IDs bleiben bestehen, während Klassen bereinigt und von Duplikaten befreit werden. Liegt neben einer BEM-Basisklasse eine spezifischere Modifier-Klasse vor, wird die 
redundante Basisklasse entfernt. Zusätzlich bleiben unter anderem Attribute zur semantischen Auszeichnung sowie ausgewählte \texttt{data-*}-Attribute erhalten, sofern diese 
nicht erkennbar der Darstellung, dem Tracking oder der JavaScript-Steuerung dienen.

Nach der Attributbereinigung werden leere Container entfernt. Für die abschließende Serialisierung wählt der Filter nach Möglichkeit ein \texttt{main}-Element als semantischen 
Wurzelbereich. Ist dieses nicht vorhanden, werden der Reihe nach \texttt{article}, \texttt{body} oder das vollständige verbleibende Dokument verwendet. Die Zwischenraumformatierung 
wird vereinheitlicht, der Inhalt jedoch nicht durch ein festes Zeichenlimit gekürzt.

Das erzeugte \texttt{PreprocessingResult} enthält neben den Metadaten und dem vorverarbeiteten HTML auch den ursprünglichen Umfang, die Anzahl entfernter Elemente, die 
Reduktionsrate und den Verarbeitungszeitpunkt. Der Wert \texttt{preprocessed\_char\_count} bezeichnet dabei den Umfang der serialisierten Kombination aus Metadaten und 
vorverarbeitetem HTML und nicht ausschließlich die Länge des HTML-Strings. Da beide Extraktionszweige auf genau diesem gespeicherten Ergebnis arbeiten, wird eine gemeinsame 
und nachträglich rekonstruierbare Eingabegrundlage geschaffen.

\subsubsection{Status- und Fehlerbehandlung}

Die technische Ausführung wird durch persistierte Pipelinekontexte dokumentiert. Für eine requestbezogene Pipeline enthält der \texttt{TransformationPipelineContext} neben 
Start- und Abschlusszeitpunkt Metadaten zum Ausführungsmodus, zur Pipeline und zum verwendeten Worker. Für jeden Filter wird ein \texttt{FilterExecutionReport} mit Filtername, 
Ausführungsstatus, Laufzeit und gegebenenfalls einer Fehlermeldung erzeugt. Batchbezogene Verarbeitungsschritte verwenden einen entsprechenden \texttt{BatchPipelineContext}, 
der dem übergeordneten Batch zugeordnet wird.

Tritt innerhalb einer requestbezogenen Pipeline eine unerwartete Ausnahme auf, wird der zugehörige Pipelinekontext abgeschlossen und dem betroffenen \texttt{TransformationRequest} 
hinzugefügt. Die Anfrage wechselt anschließend in den Zustand \texttt{FAILED}, während die übrigen Worker weitere Anfragen desselben Batches verarbeiten. Eine Ausnahme innerhalb 
der batchweiten LLM-Einreichung betrifft dagegen sämtliche zu diesem Zeitpunkt gemeinsam verarbeiteten Anfragen. Der Ausführungsfehler wird im Batchkontext gespeichert und die 
betroffenen Requests werden als fehlgeschlagen markiert.

Technische Ausnahmen sind von regulär erzeugten fachlichen Fehlerergebnissen zu unterscheiden. Beispielsweise können Schemaabweichungen oder Providerfehler als strukturierte 
Ergebnisobjekte gespeichert werden, ohne dass sie zunächst als unbehandelte Ausnahme auftreten. Auch eine nicht durchführbare Evaluation ist vom operativen Produktergebnis 
getrennt. Der \texttt{EvaluationFilter} ist deshalb so konfiguriert, dass ein Fehler dieses Schritts die nachfolgende Exportverarbeitung nicht grundsätzlich verhindert.

Eine allgemeine Retry- oder Backoff-Logik ist nicht implementiert. Ebenso erfolgt keine automatische Fortsetzung durch einen Hintergrundprozess. Die Wiederaufnahme der Verarbeitung 
nach Abschluss des externen LLM-Batches wird durch einen erneuten expliziten Aufruf ausgelöst, der auf den zuvor persistierten Zuständen und Zwischenresultaten aufbaut.

\subsection{Umsetzung der LLM-basierten Extraktion}
\subsubsection{Promptverwaltung und Promptstrategie}
\subsubsection{Attributbezogene Umsetzung}
\subsubsection{Batchanbindung und Antwortverarbeitung}

\subsection{Umsetzung des regelbasierten Referenzansatzes}
\subsubsection{Kandidatenerkennung}
\subsubsection{Kontextbewertung und Ausschluss}
\subsubsection{Plausibilitätsprüfung und Konfliktlösung}

\subsection{Validierung, Evaluationsunterstützung und Systemintegration}
\subsubsection{Schema-Validierung und Normalisierung}
\subsubsection{Technische Unterstützung der Evaluation}
\subsubsection{Nachgelagerte Integration}